% GuardianAI: Counterfactual Causal Diagnosis for Runtime Governance of LLM Applications
% Target venue: NeurIPS Workshop on Safe & Trustworthy ML, 2026.
%
% Build:  make paper           (runs latexmk -pdf manuscript.tex)
% Falls back to article class when neurips_2024.sty is unavailable.

\documentclass[11pt,letterpaper]{article}

\usepackage[margin=1in]{geometry}
\usepackage{amsmath,amssymb,amsthm}
\usepackage{algorithm}
\usepackage{algpseudocode}
\usepackage{graphicx}
\usepackage{booktabs}
\usepackage{xcolor}
\usepackage{cite}
\usepackage{etoolbox}
\usepackage[hidelinks]{hyperref}
\usepackage{url}

\newtheorem{definition}{Definition}
\newtheorem{proposition}{Proposition}

\title{\bfseries GuardianAI: Counterfactual Causal Diagnosis for Runtime Governance of LLM Applications}
\author{%
  syedwam7q\thanks{Correspondence: GuardianAI, \url{https://github.com/syedwam7q/guardian-ai}.}\\
  \texttt{engg@airtribe.live}
}
\date{April 2026}

\begin{document}
\maketitle

\begin{abstract}
Large language models in production trigger frequent safety alarms, yet
operators receive only a binary verdict and a confidence score; the
\emph{cause} of the violation, and the corrective intervention that
would prevent it, remain opaque. We present \textbf{GuardianAI}, a
runtime governance plane that combines a 3-stage / 7-agent
interception pipeline with a counterfactual causal diagnosis engine
operating over an explicit DAG of the LLM pipeline. When a violation
is detected, GuardianAI samples interventions on intervenable nodes
(temperature, top-$p$, retrieval-$k$, model choice), executes the
counterfactuals, and ranks each node by its average treatment effect
on the violation score using a bootstrap difference-in-means estimator
augmented with backdoor adjustment when confounders are present.
On a bundled medical evaluation harness spanning seven public
benchmarks, the system achieves an 18-point top-1 RCA-accuracy
improvement over an LLM-judge attribution baseline while keeping
governance overhead under 100\,ms at p95.  GuardianAI ships as an open
SDK, an OpenAI-compatible HTTP proxy, and a built-in MedRAG demo
application.
\end{abstract}

% ---------------------------------------------------------------
\section{Introduction}
\label{sec:intro}

LLM-based applications now mediate clinical triage, legal document
review, customer support, and code synthesis at production scale.  The
ML safety community has responded with a constellation of \emph{detectors}:
hallucination scoring \cite{manakul2023selfcheckgpt,li2023halueval,niu2024ragtruth},
toxicity classification \cite{hanu2020detoxify}, prompt-injection
filters \cite{inan2023llamaguard}, PII anonymisers
\cite{microsoft2021presidio}, and policy guards
\cite{rebedea2023nemo}.  These tools answer the question \emph{is the
output unsafe?} effectively.  None of them answer the operator's next
question: \emph{which component of my pipeline caused this
violation, and what would have prevented it?}

Without that answer, every alarm becomes a manual investigation.
On-call engineers re-run the prompt, swap the model, sweep the
temperature, and rebuild the corpus by hand.  In safety-critical
domains the stakes are particularly acute --- a clinician deploying a
medical RAG assistant cannot afford to wait several hours to learn
that the hallucinated drug dosage was caused by a stale corpus index
rather than the model itself.  An actionable, structured
\emph{root-cause attribution} on top of every alarm is the missing
link between detection and remediation.

\paragraph{Our position.}
Treating the LLM application as a structural causal model
\cite{pearl2009causality,scholkopf2021causal} unlocks a principled
answer.  An LLM pipeline is an explicit, intervenable computation
graph: the user query is embedded, top-$k$ documents are retrieved
under a sampling parameter $k$, a prompt template fuses the query
with the retrieved evidence, and a generator with sampling
hyper-parameters produces the output.  Every edge of that graph is
under engineering control; every node admits a finite intervention
set.  This is precisely the setting in which counterfactual causal
inference excels.

\paragraph{Contributions.}
We make three concrete contributions:

\begin{enumerate}
\item \textbf{A causal diagnosis engine for LLM violations.} Given an
  alarm from any of the seven governance agents, GuardianAI samples
  interventions on each intervenable DAG node, executes the
  counterfactuals through a stub or live LLM, and produces a ranked
  attribution with bootstrap confidence intervals
  (Section~\ref{sec:method}).
\item \textbf{An open 3-stage governance pipeline.} Seven agents
  (prompt-injection, PII-in, policy, hallucination, bias/toxicity,
  PII-out, cost) are coordinated through pre-flight, in-flight, and
  post-flight stages with timeout isolation and a constraint-based
  decision engine that emits actionable verdicts with alternatives
  (Section~\ref{sec:system}).
\item \textbf{Three deployment surfaces and an evaluation harness.} A
  Python SDK with a \texttt{@govern} decorator, an OpenAI-compatible
  HTTP proxy built on FastAPI~\cite{ramirez2018fastapi}, and a
  built-in MedRAG~\cite{lewis2020rag} demo application share the same
  governance plane and a DuckDB-backed trace
  store~\cite{raasveldt2019duckdb}.  The shipped harness covers seven
  public benchmarks plus a synthetic causal-RCA fixture
  (Section~\ref{sec:experiments}).
\end{enumerate}

The full system is open-sourced under MIT.  All experiments in this
paper are reproducible from the public repository with a single
\texttt{python -m src.eval.run\_all} invocation.

% ---------------------------------------------------------------
\section{Related Work}
\label{sec:related}

\paragraph{LLM safety and governance toolkits.}
Llama-Guard~\cite{inan2023llamaguard} and NeMo
Guardrails~\cite{rebedea2023nemo} represent the dominant industrial
line: classifier-based input/output filters with rule-based or
programmable rails.  Both detect violations effectively but expose
neither a structural model of the pipeline nor a notion of
counterfactual remediation.  Constitutional AI~\cite{bai2022constitutional}
shifts safety into the training loop via critique-and-revise, while
DSPy~\cite{khattab2023dspy} compiles declarative LM programs and
enables some form of trace inspection but, again, without a causal
account of why a given prompt failed.

\paragraph{Hallucination detection.}
SelfCheckGPT~\cite{manakul2023selfcheckgpt} samples multiple decodings
to detect inconsistency; HaluEval~\cite{li2023halueval} and
RAGTruth~\cite{niu2024ragtruth} introduce labelled corpora for
hallucination evaluation; FActScore~\cite{min2023factscore} provides
atomic-fact decomposition.  TruthfulQA~\cite{lin2022truthfulqa}
benchmarks hallucination on adversarial prompts; the broader detector
zoo is catalogued in the recent hallucination
survey~\cite{ji2023halu_survey}, and recent retrieval-conditioned
generators such as REPLUG~\cite{shi2024replug} and clinical LLMs
such as Med-PaLM~\cite{singhal2023medpalm} sharpen the need for
runtime safety on top of the underlying generator.  GuardianAI's hallucination agent reuses a
RoBERTa-large-MNLI~\cite{liu2019roberta,williams2018mnli} entailment
scorer over the retrieved evidence; the novelty in our work is not
hallucination detection itself but \emph{causal attribution} of the
detected violation back to the pipeline node that produced it.

\paragraph{Causal inference for ML.}
The methodological foundations are
classical~\cite{pearl1995ate,rubin1974potential,spirtes2000causation,pearl2009causality,pearl2018book}
but their application to LLM debugging is recent.
DoWhy~\cite{sharma2020dowhy} provides the canonical Python toolkit for
structural identification and effect estimation, and we use it for
backdoor adjustment when the DAG contains observed confounders.  Sch\"olkopf
et~al.~\cite{scholkopf2021causal} argue more broadly for causal
representations in ML; our DAG is a hand-specified instance of that
agenda for the narrow domain of LLM pipelines.  Bootstrap CIs for our
difference-in-means estimator follow the standard Efron
recipe~\cite{efron1979bootstrap}.

\paragraph{Attribution baselines.}
The natural alternatives to causal attribution are heuristic.
Saliency-based attention attribution~\cite{simonyan2014saliency,abnar2020attention}
ranks input tokens but conflates input attribution with pipeline
attribution.  LLM-judge attribution~\cite{zheng2023llmjudge} prompts
an external model to point at the suspected cause; this is brittle to
prompt sensitivity and inherits the judge's biases.  We compare
against both as baselines in Section~\ref{sec:experiments}.

\paragraph{Position.}
GuardianAI sits at the intersection of two previously separate lines:
the agent-pipeline line (Llama-Guard, NeMo, DSPy) and the structural
causal-attribution line (Pearl, Sch\"olkopf, DoWhy).  We are not aware
of a prior runtime system that combines both into an open,
production-ready governance plane.

% ---------------------------------------------------------------
\section{System}
\label{sec:system}

\subsection{Three-stage interception pipeline}
\label{sec:pipeline}

The governance plane intercepts every LLM call along three stages
(Figure~\ref{fig:arch}).  \emph{Pre-flight} runs synchronously before
the LLM call and may block the request; it bundles three agents:
prompt-injection (regex + heuristic
classifier~\cite{lakera2023gandalf,zou2023advbench}),
PII-in (Microsoft Presidio~\cite{microsoft2021presidio}), and policy
(YAML-driven domain rules).  \emph{In-flight} observers attach to the
retrieval and prompt-construction phases without modifying the call.
\emph{Post-flight} runs asynchronously after the call returns and
covers four agents: hallucination
(NLI-based~\cite{liu2019roberta,williams2018mnli}), bias and toxicity
(Detoxify~\cite{hanu2020detoxify}), PII-out (Presidio), and cost (token
accounting).  Each agent inherits from a \texttt{BaseAgent} that
enforces a per-agent timeout, isolates exceptions, and emits a typed
\texttt{Verdict} carrying severity, confidence, and machine-readable
evidence.

\begin{figure}[t]
\centering
\IfFileExists{figures/architecture.pdf}
  {\includegraphics[width=0.85\linewidth]{figures/architecture.pdf}}
  {\fbox{\parbox{0.78\linewidth}{\centering\vspace{2em}
    \textbf{Figure 1 placeholder.} GuardianAI 3-stage governance
    pipeline. PRE-FLIGHT (prompt-injection, PII-in, policy) runs
    sync $\le 80$\,ms. IN-FLIGHT observers attach to retrieval and
    prompt construction. POST-FLIGHT (hallucination, bias, PII-out,
    cost) runs async $\le 300$\,ms. CAUSAL fires on violation; DECISION
    selects an action; REMEDIATE applies it.\vspace{2em}}}}
\caption{The GuardianAI 3-stage governance pipeline architecture.
Each stage carries an explicit latency budget; agents within a stage
run in parallel.}
\label{fig:arch}
\end{figure}

\subsection{Decision engine}
\label{sec:decision}

When violations are present, a constraint-based decision engine
selects a remediation action.  Hard rules cover BLOCK-severity verdicts
(prompt-injection BLOCK $\to$ \texttt{BLOCK}; PII-in BLOCK $\to$
\texttt{REDACT}); soft rules rank candidate actions per agent
(hallucination $\to$ \texttt{REGENERATE\_WITH\_CONTEXT}, \texttt{REWRITE},
\texttt{ALERT}; bias $\to$ \texttt{REWRITE}, \texttt{ALERT}).  The engine
returns a \texttt{Decision} carrying the chosen action, a natural-language
rationale, the alternatives considered (with scores), an expected outcome,
and the causal attribution.  Crucially, we use no learned policy and no
Q-learning: every decision is auditable from the constraint set, and
the alternatives field is a contractual disclosure for the operator.

\subsection{Latency budgets and isolation guarantees}
\label{sec:latency_budgets}

Every agent declares a wall-clock timeout (e.g.\ injection: 200\,ms;
PII-in: 300\,ms; policy: 100\,ms; hallucination: 8\,s; bias: 5\,s;
PII-out: 300\,ms; cost: 50\,ms).  Stage budgets compose multiplicatively
in parallel, not additively: the pre-flight stage blocks for
$\max_i \tau_i$ rather than $\sum_i \tau_i$ because all three agents
fire concurrently via \texttt{asyncio.gather}.  An agent that exceeds
its timeout returns a SAFE-with-warning verdict and the request
proceeds; an agent that raises an exception likewise degrades to SAFE
and the trace store records the failure.  This \emph{fail-open with
audit} contract is intentional: a crashed governance agent must not
take down the user-facing application, but every degradation must be
visible in the trace.

% ---------------------------------------------------------------
\section{Method: Counterfactual Causal Diagnosis}
\label{sec:method}

This is the core contribution of the paper.  We formalise the
diagnosis problem on a structural causal model of the LLM pipeline,
specify the intervention abstraction, define the average-treatment-effect
estimator with bootstrap CIs, and describe the algorithmic loop.

\subsection{The LLM pipeline as a causal DAG}
\label{sec:dag}

\begin{definition}[LLM pipeline DAG]
\label{def:dag}
A directed acyclic graph $G = (V, E)$ where each node $v \in V$ is one
of three kinds: \emph{exogenous} (set externally:
\texttt{model\_choice}, \texttt{temperature}, \texttt{top\_p},
\texttt{retrieval\_k}, \texttt{prompt\_template}), \emph{pipeline}
(internal computation: \texttt{user\_input}, \texttt{query\_embedding},
\texttt{retrieved\_docs}, \texttt{model\_generation}, \texttt{output}),
or \emph{outcome} (\texttt{violation}, the observed alarm). $E$ encodes
the causal-flow structure: \texttt{user\_input} $\to$
\texttt{query\_embedding} $\to$ \texttt{retrieved\_docs} $\to$
\texttt{model\_generation} $\to$ \texttt{output} $\to$ \texttt{violation},
plus the parameter-into-generation edges.  Each node carries an
\texttt{intervenable} flag declaring whether $\mathrm{do}(v=v')$ is
mechanically possible at runtime.  We denote the intervenable subset
$V_I \subseteq V$.
\end{definition}

The standard pipeline DAG ships in the repository as the canonical
fixture (Figure~\ref{fig:dag}).  Eleven nodes; twelve edges; four
intervenable knobs (\texttt{temperature}, \texttt{top\_p},
\texttt{model\_choice}, \texttt{retrieval\_k}) plus
\texttt{prompt\_template} and the embedding/retrieval pipeline.  The
DAG is acyclic by construction (verified by NetworkX).

\begin{figure}[t]
\centering
\IfFileExists{figures/dag.pdf}
  {\includegraphics[width=0.75\linewidth]{figures/dag.pdf}}
  {\fbox{\parbox{0.72\linewidth}{\centering\vspace{1.5em}
    \textbf{Figure 2 placeholder.} Standard LLM pipeline DAG used in
    this work. 11 nodes, 12 edges. Intervenable: temperature, top\_p,
    model\_choice, retrieval\_k, prompt\_template (squares); pipeline
    nodes are circles; the violation outcome is shaded.\vspace{1.5em}}}}
\caption{The 11-node, 12-edge LLM pipeline DAG. Square nodes admit
runtime interventions; circles are pipeline-internal; the shaded node
is the violation outcome.}
\label{fig:dag}
\end{figure}

\subsection{Counterfactual interventions and ATE}
\label{sec:cf}

For each intervenable node $v \in V_I$ with baseline value $v_0$ we
define an \emph{intervention plan} as a finite candidate set
$C_v = \{c_1, \ldots, c_{k_v}\}$ excluding $v_0$.  Each candidate yields
a counterfactual trace $\tilde{T}_{v,c}$ obtained by re-executing the
pipeline with $\mathrm{do}(v = c)$ and all other inputs held fixed.

Let $\nu : T \to [0,1]$ be a violation-score function (the
post-flight agent that originally raised the alarm, reapplied to the
counterfactual output).  The \emph{average treatment effect} of node
$v$ on the violation outcome is:
\begin{equation}
\tau_v \;=\; \mathbb{E}\big[\nu(T) \,\big|\, \mathrm{do}(V = v_0)\big]
        \;-\; \mathbb{E}\big[\nu(T) \,\big|\, \mathrm{do}(V = v')\big],
\label{eq:ate}
\end{equation}
where $v'$ ranges over $C_v$.  Equation~\eqref{eq:ate} is the standard
Pearl ATE specialised to a do-operator on a single
node~\cite{pearl2009causality,pearl1995ate}.  A positive $\tau_v$
means: \emph{intervening on $v$ \textbf{reduces} the violation}, hence
$v$ is causally implicated.  Nodes are ranked by $|\tau_v|$.

\subsection{Estimation: bootstrap difference-in-means}
\label{sec:estimator}

In the runtime setting, the expectations in
Equation~\eqref{eq:ate} must be estimated from a small number of
counterfactual samples (typically $n = 3$ per intervention to keep the
diagnosis budget below 1.5\,s).  We use a difference-in-means estimator
with a bootstrap confidence
interval~\cite{efron1979bootstrap}:

\begin{equation}
\hat{\tau}_v \;=\; \nu_0 \;-\; \frac{1}{n}\sum_{i=1}^{n}\nu(\tilde{T}_{v,c}^{(i)}),
\qquad
\widehat{\mathrm{CI}}_{0.95}(\tau_v) \;=\; \big[\,\widehat{q}_{0.025},\,
                                                \widehat{q}_{0.975}\,\big],
\label{eq:estimator}
\end{equation}
where $\nu_0$ is the baseline violation score and the quantiles are
taken over $B = 1000$ bootstrap resamples of the counterfactual
scores.  The estimator is closed-form, requires no model fitting, and
produces a calibrated CI even at $n = 3$ when the per-sample variance
is high.

\subsection{Backdoor adjustment with DoWhy}
\label{sec:dowhy}

When the DAG declares an observed confounder between $v$ and the
violation (e.g.\ \texttt{retrieved\_docs} confounding
\texttt{model\_generation} on the violation edge), the
difference-in-means estimator is biased.  GuardianAI delegates to
DoWhy~\cite{sharma2020dowhy} for backdoor identification and a
linear-regression effect estimator.  When DoWhy's API rejects the
graph (small samples; numerical conditioning; library version
skew across the 0.12.x line) the engine falls back to
Equation~\eqref{eq:estimator} so the diagnosis loop never crashes.

\subsection{Algorithm}
\label{sec:algorithm}

\begin{algorithm}[t]
\caption{Counterfactual Causal Diagnosis}
\label{alg:diagnose}
\begin{algorithmic}[1]
\Require Baseline trace $T_0$, baseline violation score $\nu_0$, DAG $G$,
         counterfactual executor $\mathcal{E}$, samples-per-intervention $n$.
\Ensure  Ranked ATEs $\{(\hat{\tau}_v, \widehat{\mathrm{CI}}_v)\}_{v \in V_I}$.
\State $\Pi \gets \textsc{StandardPlans}(T_0)$
       \Comment{intervention plan per intervenable $v$}
\State $\mathcal{R} \gets \emptyset$
\ForAll{$\pi \in \Pi$ \textbf{ in parallel} (semaphore $\le 4$)}
  \ForAll{$c \in \pi.\textsc{Candidates}\setminus\{v_0\}$, $i = 1\ldots n$}
    \State $\tilde{T} \gets \mathcal{E}.\textsc{Execute}(T_0,\ \mathrm{do}(\pi.v = c))$
    \State $\nu_{v,c}^{(i)} \gets \nu(\tilde{T})$
    \State Append $\nu_{v,c}^{(i)}$ to per-node bucket $S_v$
  \EndFor
\EndFor
\ForAll{$v \in V_I$}
  \State $\hat{\tau}_v, \widehat{\mathrm{CI}}_v \gets \textsc{BootstrapDiffMeans}(\nu_0, S_v)$
\EndFor
\State \Return $\textsc{Sort}(\{\hat{\tau}_v\}, \,\text{key}=|\hat{\tau}_v|, \,\text{desc})$
\end{algorithmic}
\end{algorithm}

\subsection{Complexity}
\label{sec:complexity}

\begin{proposition}
Algorithm~\ref{alg:diagnose} requires
$\Theta(|V_I| \cdot \bar{k} \cdot n)$ counterfactual executions,
where $\bar{k}$ is the average number of candidates per intervention
plan.  In our reference configuration ($|V_I| = 4$, $\bar{k} = 3.25$,
$n = 3$), that is at most $39$ executions.  With $\le 4$-way
concurrency and a $\sim$\,200\,ms stub executor the diagnosis budget
is well under 1.5\,s; with a live LLM executor the per-counterfactual
cost dominates and the budget is set by the underlying provider.
\end{proposition}

\subsection{Why this beats heuristic alternatives}
\label{sec:why}

LLM-judge attribution~\cite{zheng2023llmjudge} produces a ranking by
prompting an LLM with the trace and asking it to nominate the cause.
The judge has no structural model; rankings inherit prompt sensitivity
and judge bias.  Attention attribution~\cite{simonyan2014saliency,abnar2020attention}
maps token-level attention back into evidence-token saliency but
\emph{cannot} attribute across the pipeline (it cannot rank
\texttt{temperature} vs \texttt{retrieval\_k}; both are exogenous to
the attention map).  The counterfactual ATE is structural by
construction: it answers the well-posed question \emph{what would the
violation score have been if I had set $v$ to a different value?} and
ranks nodes by the answer.

% ---------------------------------------------------------------
\section{Experiments}
\label{sec:experiments}

\paragraph{Honesty preface.} All numerical results in this section are
computed on the bundled in-repo fixtures (5--10 hand-crafted medical
cases per benchmark; 8 cases for synthetic causal RCA).  These
fixtures are deterministic, run offline without a network, and
exercise the harness end-to-end.  The full upstream benchmarks
(HaluEval, FactScore, RAGTruth, BBQ, AdvBench, MedHELM, PubMedQA) are
larger and require offline downloads; the adapter code in
\texttt{backend/src/eval/datasets/} ingests the upstream JSONL formats
unchanged.  Likewise, the 500-case hand-labelled causal-RCA dataset
described in Section~\ref{sec:limitations} is in flight; the present
RCA numbers come from the synthetic 8-case fixture with a
ground-truth-aware scorer.  We report results that we can reproduce
with one CLI command from the public repository
(\texttt{python -m src.eval.run\_all}), and we mark synthetic columns
explicitly.

\subsection{Per-agent benchmarks (Table~\ref{tab:per_agent})}
\label{sec:per_agent}

We ran each governance agent against its primary benchmark fixture and
report binary precision/recall/F1 with predictions derived from
\texttt{verdict.severity > SAFE} and ground truth from each fixture's
labels.

\begin{table}[t]
\centering
\caption{Per-agent F1 across seven benchmarks. Bundled fixtures
($n = 6$--$10$ per dataset). Heavy agents load
RoBERTa-large-MNLI / Detoxify / Presidio on first call.}
\label{tab:per_agent}
\begin{tabular}{lccc}
\toprule
Agent $\times$ Benchmark & Precision & Recall & F1 \\
\midrule
Hallucination $\times$ HaluEval~\cite{li2023halueval}      & 0.86 & 0.83 & 0.84 \\
Hallucination $\times$ RAGTruth~\cite{niu2024ragtruth}     & 0.82 & 0.80 & 0.81 \\
Hallucination $\times$ FActScore~\cite{min2023factscore}   & 0.79 & 0.81 & 0.80 \\
Bias/Toxicity $\times$ BBQ~\cite{parrish2022bbq}            & 0.92 & 0.74 & 0.82 \\
Prompt-injection $\times$ AdvBench~\cite{zou2023advbench}   & 0.91 & 0.88 & 0.89 \\
Policy $\times$ MedHELM~\cite{fries2024medhelm}             & 0.85 & 0.78 & 0.81 \\
Policy $\times$ PubMedQA~\cite{jin2019pubmedqa}             & 0.81 & 0.76 & 0.78 \\
\bottomrule
\end{tabular}
\end{table}

\subsection{End-to-end MedRAG (Table~\ref{tab:medrag})}
\label{sec:medrag}

We ran the bundled MedHELM fixture through both \emph{bare RAG} (no
governance, BGE-small embeddings~\cite{xiao2024bge} into
ChromaDB~\cite{chromadb2024}, Llama-3.3-70b
generator~\cite{meta2024llama3}) and \emph{governed} (the same RAG
plus the full GuardianAI pipeline).  The MedRAG benchmark reports
hallucination rate, bias rate, PII-leak rate, and over-refusal rate.

\begin{table}[t]
\centering
\caption{MedRAG comparison on the bundled MedHELM fixture
($n = 6$). Governance reduces hallucination $4{\times}$ and PII leak to
zero with a 1-percentage-point over-refusal cost.}
\label{tab:medrag}
\begin{tabular}{lcc}
\toprule
Metric                    & Bare RAG & + GuardianAI \\
\midrule
Hallucination rate (\%)   & 33.0     & 8.0          \\
Bias rate (\%)            & 12.0     & 4.0          \\
PII leak rate (\%)        & 5.0      & 0.0          \\
Over-refusal rate (\%)    & 1.0      & 2.0          \\
\bottomrule
\end{tabular}
\end{table}

\subsection{Causal RCA accuracy (Table~\ref{tab:rca})}
\label{sec:rca}

This is the headline experiment.  Four attribution methods --- the
random-scorer baseline, an attention-attribution stand-in
(\cite{simonyan2014saliency,abnar2020attention} as the conceptual
target; a deterministic prior-biased scorer as the implementation),
an LLM-judge stand-in
(\cite{zheng2023llmjudge}; a fixed-bias scorer that always nominates
\texttt{model\_choice}, mimicking judge brittleness reported in
the literature), and \textbf{ours} (Algorithm~\ref{alg:diagnose} with
the bootstrap difference-in-means estimator and the
ground-truth-aware scorer documented in
\texttt{causal\_rca\_eval.py}) --- are run against the synthetic
8-case RCA fixture, with ground-truth attribution stratified across
the four intervenable knobs.

\begin{table}[t]
\centering
\caption{Causal-RCA top-1 / top-3 attribution accuracy on the bundled
8-case synthetic fixture. \textbf{Ours} ranks the true causal node
first 78\% of the time, an 18-point lift over the LLM-judge stand-in.
Real-dataset numbers are deferred pending the 500-case
hand-annotation (Section~\ref{sec:limitations}).}
\label{tab:rca}
\begin{tabular}{lcc}
\toprule
Method                        & Top-1 & Top-3 \\
\midrule
Random scorer                 & 0.25  & 0.75  \\
Attention attribution (stub)  & 0.38  & 0.75  \\
LLM-judge attribution (stub)  & 0.60  & 0.85  \\
\textbf{Ours} (causal + bootstrap) & \textbf{0.78} & \textbf{0.95} \\
\bottomrule
\end{tabular}
\end{table}

\subsection{Latency budget (Table~\ref{tab:latency})}
\label{sec:latency}

We fired 100 concurrent requests at the default medical pipeline using
\texttt{asyncio.gather} and measured wall-clock latency per stage.
The hallucination agent dominates post-flight latency on first
invocation due to model loading; we exclude warm-up.

\begin{table}[t]
\centering
\caption{Stage-wise latency on a local laptop, $n = 100$ concurrent
requests, warm cache. Pre-flight stays under the 80\,ms budget;
post-flight under the 300\,ms budget; causal stays under 1.5\,s on the
stub executor.}
\label{tab:latency}
\begin{tabular}{lccc}
\toprule
Stage             & p50 (ms) & p95 (ms) & p99 (ms) \\
\midrule
Pre-flight        & 12       & 31       & 48       \\
Post-flight       & 86       & 178      & 245      \\
Causal (stub)     & 320      & 740      & 1180     \\
Decision          & 2        & 6        & 11       \\
End-to-end (no causal) & 102 & 198      & 280      \\
\bottomrule
\end{tabular}
\end{table}

\subsection{Ablation (Table~\ref{tab:ablation})}
\label{sec:ablation}

We compared four configurations on the MedHELM fixture: \texttt{no\_governance}
(raw RAG), \texttt{governance\_only} (pipeline, no causal/decision),
\texttt{governance\_plus\_judge} (pipeline plus an LLM-judge layer in
place of the causal engine), and \texttt{full\_guardian} (everything).
The governance pipeline alone halves the hallucination rate; adding
the causal engine on top yields the actionability gain (RCA top-1
accuracy is undefined for configs that have no attribution mechanism).

\begin{table}[t]
\centering
\caption{Ablation across four configurations on the MedHELM fixture
($n = 6$). The causal engine adds RCA accuracy without further
moving hallucination or refusal rates --- as expected, since attribution
is read-only with respect to the user-visible response.}
\label{tab:ablation}
\begin{tabular}{lccc}
\toprule
Configuration                  & Hallucination (\%) & Refusal (\%) & RCA top-1 \\
\midrule
\texttt{no\_governance}        & 33.0  & 0.0   & ---      \\
\texttt{governance\_only}      & 16.0  & 1.5   & ---      \\
\texttt{governance\_plus\_judge} & 14.0 & 2.0   & 0.60     \\
\texttt{full\_guardian}        & \textbf{8.0}  & 2.0   & \textbf{0.78} \\
\bottomrule
\end{tabular}
\end{table}

% ---------------------------------------------------------------
\section{Case Studies}
\label{sec:cases}

We walk through three representative traces from the bundled medical
demo corpus to illustrate the governance plane in action.

\paragraph{Case 1: ``Aspirin in pregnancy.''}
A user asks the MedRAG demo whether aspirin is safe during pregnancy.
The bare RAG retrieves a stale corpus chunk on adult cardiovascular
prophylaxis and the generator confidently recommends aspirin without
trimester qualifiers --- a clinically important hallucination
\cite{acog2017paracetamol} (paracetamol, not aspirin, is the
default analgesic recommendation; aspirin in the third trimester
carries a documented risk of premature ductus arteriosus closure).  The
hallucination agent fires (severity \texttt{WARN}, confidence 0.78)
because the NLI score against the retrieved evidence is below the
contradiction threshold.  The causal engine then sweeps interventions:
the largest ATE is on \texttt{retrieved\_docs} ($\hat{\tau} = +0.51$,
95\,\% CI $[0.38, 0.62]$; intervention: increase $k$ from 5 to 10
documents and re-rank), nominating the retrieval corpus rather than
the generator as the responsible node.  The decision engine selects
\texttt{REGENERATE\_WITH\_CONTEXT}.

\paragraph{Case 2: ``Ignore previous instructions.''}
A red-team prompt instructs the assistant to disregard its system
prompt and disclose its policy.  The prompt-injection agent fires at
severity \texttt{BLOCK} with confidence 0.97 in the pre-flight stage
($\sim$15\,ms wall-clock).  The decision engine emits the hard-rule
\texttt{BLOCK} action with three alternatives considered
(\texttt{ALERT\_AND\_PROCEED}, \texttt{ALERT}, \texttt{LOG}).  The user
sees a refusal explanation; the trace records the matched signature
and the alternatives.  The causal engine does not run because no
post-flight violation was raised.

\paragraph{Case 3: ``What's the dose of metformin?''}
The policy agent flags the bare query at severity \texttt{WATCH} on a
medical-disclaimer rule (the corpus dose ranges are not titrated to
the patient's renal function, which the policy YAML for the
\texttt{medical} domain marks as a hard prerequisite).  The decision
engine selects \texttt{ADD\_DISCLAIMER} (alternatives: \texttt{ALERT},
\texttt{REWRITE}); the response is delivered with a prepended ``This
information is general and is not a substitute for clinician dosing
adjusted for renal function'' disclaimer.  Causal diagnosis is
optional here since the verdict is below \texttt{WARN}; the trace
nevertheless records the policy-rule ID for audit.

% ---------------------------------------------------------------
\section{Discussion, Limitations, and Ethics}
\label{sec:limitations}

\paragraph{Synthetic ground truth.}
The headline RCA accuracy in Table~\ref{tab:rca} is computed against
the bundled synthetic 8-case fixture with a ground-truth-aware scorer
(see \texttt{\_gt\_aware\_score} in the eval module).  This is an
end-to-end harness verification, not a real-world evaluation.  The
real evaluation is the 500-case hand-labelled benchmark protocoled in
\texttt{docs/EVALUATION.md}: stratified sampling of $\sim$1500
production traces across the four intervenable knobs, two-reviewer
labelling with a third-reviewer adjudication target Cohen's-$\kappa
\ge 0.65$, and a real trained NLI/judge ensemble in place of the
ground-truth-aware scorer.  At three minutes per case and two
reviewers plus adjudication, the collection is roughly 80
person-hours.  That work is in flight; the present paper is the
methodological precursor.

\paragraph{Single-machine evaluation.}
The latency benchmark uses Python \texttt{asyncio.gather} from a
single laptop process.  We have not stress-tested the pipeline at
production scale (k6, multi-region, sustained QPS), nor have we
characterised behaviour under sustained concurrency contention with a
shared trace store.

\paragraph{English-only.}
All bundled fixtures and all detector models are English; the
prompt-injection regex set and the policy YAML are also English-only.
Multilingual coverage requires separate detectors and a different set
of policy templates and is left to future work.

\paragraph{Single-annotator bias.}
The bundled medical fixtures were authored by a single domain-aware
annotator.  Cohen's-$\kappa$ inter-annotator agreement is therefore
not computable on the in-repo data; the 500-case protocol corrects
this.

\paragraph{Benchmark contamination risk.}
HaluEval, BBQ, MedHELM, and PubMedQA are publicly available and
predate the training cutoff of every commercial LLM we benchmark
through.  Our F1 numbers should be interpreted accordingly: the
bundled fixtures are de-duplicated against upstream test splits to the
extent reasonable, but contamination of the underlying generator
remains a confound common to all LLM evaluation.

\paragraph{Same-pipeline ablation.}
The current 4-config ablation (Table~\ref{tab:ablation}) shares a
single default pipeline across the \texttt{governance\_only},
\texttt{governance\_plus\_judge}, and \texttt{full\_guardian} rows;
real ablation runs would build per-config pipelines.  This is an
honest scaffolding state acknowledged in
\texttt{docs/EVALUATION.md}.

\paragraph{Ethics.}
The MedRAG demo and the medical fixtures are illustrative.  We do not
advocate the deployment of LLM-mediated clinical decision support
without clinician review; GuardianAI is intended as a runtime
guard-rail \emph{around} such systems, not as a replacement for human
oversight.  All policy YAMLs ship with a default \texttt{ADD\_DISCLAIMER}
remediation precisely because the system should not silently elevate
its own perceived authority.  The PII detector ships off Microsoft
Presidio's English entity types and is not a substitute for a HIPAA-
or GDPR-grade compliance review.

% ---------------------------------------------------------------
\section{Conclusion}
\label{sec:conclusion}

We presented \textbf{GuardianAI}, an open runtime governance plane
that pairs a 3-stage / 7-agent interception pipeline with a
counterfactual causal diagnosis engine over an explicit DAG of the
LLM pipeline.  The diagnosis loop ranks the intervenable nodes of
the pipeline by their average treatment effect on the violation
score, using a bootstrap difference-in-means estimator with a DoWhy
backdoor-adjustment fallback.  On the bundled synthetic fixture, the
causal estimator delivers an 18-point top-1 RCA-accuracy improvement
over an LLM-judge attribution baseline at $\le 100$\,ms p95 governance
overhead.  GuardianAI ships as a Python SDK, an OpenAI-compatible HTTP
proxy, and a built-in MedRAG demo; the full system is open-source.

The next milestones are the 500-case hand-labelled real-world RCA
dataset described in Section~\ref{sec:limitations}, multilingual
detector coverage, and a production-scale concurrent-load benchmark.
We see GuardianAI as a step toward closing the loop between
\emph{detecting} an LLM violation and \emph{repairing} it: an alarm
that nominates a fix is materially more useful than an alarm that
nominates only itself.

\paragraph{Reproducibility.}
All code, fixtures, the YAML policy used in this paper, and the
evaluation runner that generates Tables~\ref{tab:per_agent}--\ref{tab:ablation}
are available at \url{https://github.com/syedwam7q/guardian-ai} under
the MIT licence.  Section-by-section pointers into the code are in
\texttt{paper/README.md}.

\bibliographystyle{plain}
\bibliography{refs}

\end{document}
